\documentclass[9pt,a4paper]{article}
\usepackage[margin=1in]{geometry}
\usepackage{amsmath,amssymb}
\setlength{\parskip}{3pt}
\setlength{\parindent}{0pt}

\title{\textbf{Q2: Route Blocking for Reducing Forest Fire Spread}\\
COL761 Assignment 2 (2026) -- BITBLASTING}
\author{}
\date{}

\begin{document}
\maketitle
\vspace{-1.2em}

\textbf{Submitted files for Q2:} \texttt{A2/q2/forest\_fire.sh}, \texttt{A2/q2/report\_q2.pdf}.\\
\textbf{Implementation used:} \texttt{A2/q2/forest\_fire.cpp} (single-process, single-threaded C++).

\section*{Task 1: Optimization Formulation}
We build on the assignment's definition of the directed probabilistic graph
$G=(V,E,p)$, seed set $A_0$, and Independent-Cascade spreading process.
As clarified on Piazza, Task 1 is formulated for the standard model
(without hop constraint).

For each edge $e\in E$, define binary decision variable:
\[
 x_e =
 \begin{cases}
 1,&\text{if edge }e\text{ is blocked},\\
 0,&\text{otherwise.}
 \end{cases}
\]
Let $R(x)=\{e\in E:x_e=1\}$. The optimization problem is:
\[
\min_{x\in\{0,1\}^{|E|}}\; \sigma(R(x))
\quad\text{s.t.}\quad
\sum_{e\in E} x_e = k.
\]
Here $\sigma(R)=\mathbb{E}[|A_\infty|]$ is the expected total burned nodes on
$G_R=(V,E\setminus R,p)$.

(Optional extension used in implementation) for hop-limit $h\ge 1$, spread is
restricted to nodes with graph distance at most $h$ from $A_0$.

\section*{Task 2: NP-Hardness}
\textbf{Claim:} The route-blocking problem is NP-hard.

\textbf{Reduction source:} \emph{Minimum $b$-Union} (also called Small-Set-Union)
decision problem: given universe $U$, sets $\mathcal{S}=\{S_1,\dots,S_m\}$,
integer $b$, and threshold $L$, decide whether there exists a subfamily of size
$b$ whose union size is at most $L$. This problem is NP-hard.

\textbf{Construction.}
Given $(U,\mathcal{S},b,L)$, create a deterministic instance ($p=1$ for all edges):
\begin{itemize}
\item One seed node $s$.
\item One set-node $a_i$ for each set $S_i$.
\item One element-node $u_j$ for each element in $U$.
\item For each $i$, add edge $(s,a_i)$.
\item For each $u_j\in S_i$, add edge $(a_i,u_j)$.
\item For each set-node $a_i$, add $M$ private leaf nodes
      $p_{i,1},\dots,p_{i,M}$ with edges $(a_i,p_{i,r})$, where
      $M>|U|+m$.
\end{itemize}
Set budget $k=m-b$.

Because $M$ is very large, replacing any blocked non-source edge by blocking its
corresponding source edge $(s,a_i)$ never increases spread and strictly improves
(or preserves) objective. Hence an optimum can be taken to block only source
edges. Blocking exactly $k=m-b$ source edges leaves exactly $b$ active set-nodes,
call this family $\mathcal{C}$.

Burned count is then deterministic and equals:
\[
|A_\infty| = 1 + b + bM + \left|\bigcup_{S_i\in\mathcal{C}} S_i\right|.
\]
Therefore, there exists a size-$b$ family with union size $\le L$ iff there
exists a blocking set $R$ of size $k$ with
\[
\sigma(R) = |A_\infty| \le 1+b+bM+L.
\]
The reduction is polynomial, so route-blocking is NP-hard.

\section*{Task 3: Structural Properties of $\sigma(R)$}
\subsection*{3.1 Monotonicity}
If $R_1\subseteq R_2$, then $E\setminus R_2\subseteq E\setminus R_1$.
For every fixed live-edge realization, reachable burned nodes under $R_2$ are a
subset of those under $R_1$. Taking expectation gives
\[
\sigma(R_2)\le \sigma(R_1).
\]
So $\sigma$ is monotone non-increasing in blocked-edge set size.

\subsection*{3.2 Submodularity}
For a fixed realization $\omega$, let $f_\omega(R)$ be burned-node count after
removing $R$. For each node $v$, define indicator
$\mathbf{1}_\omega(v\text{ reachable from }A_0\text{ in }E_\omega\setminus R)$;
then $f_\omega(R)$ is the sum of these indicators.

Reachability indicators under edge deletions are submodular in $R$ (equivalently,
node-disconnection indicators are supermodular), hence $f_\omega$ is submodular.
Expectation preserves submodularity, so
\[
\sigma(R)=\mathbb{E}_\omega[f_\omega(R)]
\]
is submodular.

Thus for Task 3: $\sigma(R)$ is monotone non-increasing and submodular.

\section*{Task 4: Algorithm Design}
\subsection*{4.1 Final Algorithm}
We use a scalable \textbf{two-stage hybrid Monte-Carlo greedy heuristic}:
\begin{itemize}
\item \textbf{Reachability filtering:} ignore currently unblocked edges whose
tail is unreachable from seeds (or beyond hop limit).
\item \textbf{Stage-1 coarse scoring:} run multiple IC simulations and assign
credit backward from burned nodes to successful igniting edges.
\item \textbf{Stage-2 fine reranking:} shortlist top-$L$ edges from Stage-1,
sample common random worlds, and pick edge minimizing average spread.
\item \textbf{Timeout safety:} write output after each selected edge and also
at periodic checkpoints.
\end{itemize}

\subsection*{4.2 Pseudocode}
\begin{verbatim}
Input: G=(V,E,p), seeds S, budget k, n_random_instances r, hops h
R <- empty
write_output(R)

for t = 1..k:
    C <- unblocked and currently reachable candidate edges
    if C empty: pick any remaining edge, add to R, write_output(R), continue

    // Stage-1: coarse Monte-Carlo credit scoring
    score[e] <- 0 for e in C
    run T_c simulations of IC process on current graph:
        log successful parent edges for each newly burned node
        back-propagate node credit to igniting edges
        accumulate score[e]

    L <- top shortlist_size edges by score

    // Stage-2: fine reranking using common sampled worlds
    for each e in L: avgSpread[e] <- 0
    run T_f sampled live-edge worlds:
        for each e in L:
            compute deterministic spread if e is additionally blocked
            accumulate avgSpread[e]

    choose e* = argmin avgSpread[e]
    R <- R U {e*}
    write_output(R)

if |R| < k: fill remaining with arbitrary unblocked edges, writing each step
return first k edges in R
\end{verbatim}

\subsection*{4.3 Complexity}
Let $n=|V|$, $m=|E|$, budget $k$, coarse simulations/round $T_c$, fine worlds/round
$T_f$, shortlist size $L$ (constant in implementation, max 24).

Per round cost: $O(n+m)$ (filtering) $+\;O(T_c(n+m))$ (coarse MC)
$+\;O(T_fL(n+m))$ (fine rerank).
Total:
\[
O\Big(k(n+m)\big(1+T_c+T_fL\big)\Big).
\]
Memory is $O(n+m)$ with reusable buffers.

\subsection*{4.4 Approximation Guarantee Discussion}
Although $\sigma$ is submodular, our implemented solver is not exact global
optimization of $\sigma$; it uses Monte-Carlo estimation, proxy scoring, and
shortlisting. Therefore this implementation has \textbf{no formal constant-factor
approximation guarantee}. Also, the standard $(1-1/e)$ theorem applies to
maximization of monotone submodular objectives, not directly to this cardinality-
constrained minimization setting.

\section*{Task 5: Competitive Evaluation}
\subsection*{5.1 Setup}
Evaluation used the provided scripts:
\texttt{A2/q2/Eval/evaluate.py} and \texttt{evaluate.sh}.

Public settings:
\begin{itemize}
\item Dataset 1: $k=50$, \texttt{hops=-1}, \texttt{num\_sim=50}
\item Dataset 2: $k=20$, \texttt{hops=3}, \texttt{num\_sim=50}
\end{itemize}

\subsection*{5.2 Results (official num\_sim=50)}
\begin{center}
\begin{tabular}{|l|c|c|}
\hline
Method & Dataset 1 $\rho(R)$ & Dataset 2 $\rho(R)$ \\
\hline
Original baseline (older \texttt{forest\_fire\_baseline.cpp}) & 0.814223 & 0.877672 \\
Advanced & 0.823043 & 0.880727 \\
RIS & 0.836273 & 0.876386 \\
\textbf{Final submitted (current \texttt{forest\_fire.cpp})} & \textbf{0.834620} & \textbf{0.880727} \\
\hline
\end{tabular}
\end{center}

Final submitted output quality is balanced across both models and robust in
runtime.

\subsection*{5.3 Stability check (num\_sim=500)}
\begin{center}
\begin{tabular}{|l|c|c|}
\hline
Method & Dataset 1 $\rho(R)$ & Dataset 2 $\rho(R)$ \\
\hline
Original baseline & 0.746929 & 0.880560 \\
Advanced & 0.779893 & 0.884162 \\
RIS & 0.769905 & 0.880090 \\
\textbf{Final submitted} & \textbf{0.795836} & \textbf{0.884162} \\
\hline
\end{tabular}
\end{center}

\subsection*{5.4 Timeout compliance}
The solver periodically writes selected edges and also writes after each chosen
edge. Hence, in case of timeout, partial output remains valid for partial
credit.

\end{document}
